% Solution: Milgrom & Roberts (1992) Ch. 8 Mathematical Exercise 4. Input by \solitem. Check: check_T1.py (item 5).
\begin{exsol}{Milgrom \& Roberts Ch.~8, Mathematical Exercise 4}{286}{Foundations}{2}
\asked In the trust game of Ch.~8 (offeror, decision maker), the parties meet exactly twice, then exactly
three times, then any known finite number of times. Would trust be offered and honoured? Is the prediction
credible?

\method Backward induction from the last round. The stage payoffs of Ch.~8 (pp.~261--263) are: offer and
honour $(V, V)$; offer and not honour $(-L, V + G)$; no offer $(0, 0)$, with $V, G, L > 0$. In one
play, a self-interested decision maker does not honour ($V + G > V$), so the offeror does not offer
($-L < 0$).

\solution
\textbf{Two interactions.} In round~2 (final), nothing follows, so the one-shot logic applies: trust is not
honoured if offered, and so it is not offered. Round-2 payoffs are $(0,0)$ whatever happened in round~1.
In round~1, honouring cannot earn a reputation, since round~2's outcome does not depend on round~1
(the hint's question). The decision maker therefore gains $G$ by not honouring, so he will not honour,
and the offeror, anticipating this, does not offer. Both parties get $0$ in each round.

\textbf{Three interactions and any finite $N$.} By induction: suppose rounds $k+1, \dots, N$ all give
$(0,0)$ regardless of history. Then round $k$ is effectively one-shot, with the same conclusion. The
unravelling reaches round~1, so for every known finite $N$ trust is never offered and, if it were,
never honoured. The formal SPE is unique.

\textbf{Credibility.} The prediction is not credible as a description of behaviour. Experiments with
finitely repeated trust and social-dilemma games show cooperation for most of the horizon, with collapse
only near the end. Explanations: (i) people are unsure of the horizon (a probability of continuation
restores trust, item~16); (ii) some players are honest by type, so a selfish player gains by mimicking
them (reputation, Lectures~7--8, e.g.\ Kreps--Milgrom--Roberts--Wilson); (iii) a small departure from
rationality ($\varepsilon$-equilibrium) sustains cooperation for long horizons.

\interp The finite-horizon paradox shows that trust rests on the \emph{shadow of the future}: repetition
disciplines behaviour only if the interaction has no known last round.

\verify \texttt{check\_T1.py} (item~5) runs backward induction on the repeated trust game for $N = 1, 2,
3, 10$ (with $V = 3$, $G = 2$, $L = 5$): in every round the decision maker dishonours and the offeror does
not offer, giving continuation $(0,0)$. The exercise is discursive and has no printed answer.

\answer{Round $N$ is one-shot, so trust is neither offered nor honoured; honouring in round $N - 1$
cannot buy a future reward; by induction the same holds in every round. For any known finite horizon no
trust is offered. This is not credible empirically: with an uncertain end date, honest types or small
departures from rationality, cooperation reappears.}
\end{exsol}
